% \documentclass{article}
% \usepackage{amsmath}
% \usepackage{amssymb}
% \usepackage{hyperref}

\subsection{Synchrotron Undulator Emission: a summary based on \cite{SanchezdelRioReyesHerrera2025}}
% \author{Based on: Modelling undulators in ray tracing simulations}
% \date{}

% \begin{document}

% \maketitle

\subsubsection{Basic Principles}

Undulators produce brilliant, collimated synchrotron radiation through periodic electron trajectories in a sinusoidal magnetic field.

\begin{equation}
K \approx 93.37 \cdot B[T] \cdot \lambda_u[m]
\end{equation}
where \(K\) is the deflection parameter, \(B\) is the magnetic field, and \(\lambda_u\) is the undulator period.

The resonance wavelength (on-axis, \(\theta=0\)) is:
\begin{equation}
\lambda = \frac{1 + K^2/2}{2\gamma^2} \lambda_u
\end{equation}
where \(\gamma \approx 1957 \cdot E[GeV]\) is the Lorentz factor.

\subsubsection{Angular Distribution (Divergence)}

At resonance, the far-field intensity follows a sinc\(^2\) pattern:
\begin{equation}
I(\theta) \propto \text{sinc}^2\left( \frac{\pi}{2} \tilde{\theta}^2 \right), \quad \tilde{\theta} = \theta \sqrt{L/\lambda}
\end{equation}
where \(L = N\lambda_u\) is the undulator length.

The Gaussian approximation for divergence (r.m.s.) is:
\begin{equation}\label{eq:elleaumeDivergence}
\sigma_{r'} \approx 0.69 \sqrt{\frac{\lambda}{L}} \quad \text{(Elleaume model)}
\end{equation}

\subsubsection{Source Size (Spatial Distribution)}

Obtained by backpropagating far-field radiation to the source plane (\(z=0\)):
\begin{equation}\label{eq:elleaumeSize}
\sigma_r \approx 0.218 \sqrt{\lambda L}
\end{equation}

The phase space product (non-Gaussian) is:
\begin{equation}
\sigma_r \sigma_{r'} \approx \frac{\lambda}{2\pi} \quad (\text{approximately twice the diffraction limit})
\end{equation}

\subsubsection{Electron Beam Effects}

The electron beam is described by a 5D Gaussian distribution:
\begin{equation}
(x, x', y, y', \delta_{\mathcal{E}})
\end{equation}
with sizes \(\sigma_x, \sigma_y\), divergences \(\sigma_{x'}, \sigma_{y'}\), and energy spread \(\delta_{\mathcal{E}}\).

The photon source size and divergence including electron effects:
\begin{equation}\label{eq:convNoDispersion}
\Sigma_x = \sqrt{\sigma_r^2 + \sigma_x^2}, \quad \Sigma_{x'} = \sqrt{\sigma_{r'}^2 + \sigma_{x'}^2}
\end{equation}

Energy spread correction (Tanaka model):
\begin{equation}\label{eq:convDispersion}
\Sigma_{x'} = \sqrt{(Q_a \sigma_{r'})^2 + \sigma_{x'}^2}, \quad \Sigma_x = \sqrt{(Q_s \sigma_r)^2 + \sigma_x^2}
\end{equation}
where \(Q_a\) and \(Q_s\) are corrective functions depending on:
\begin{equation}
\sigma_e = 2\pi n N \delta_{\mathcal{E}}
\end{equation}
with \(n\) = harmonic number and \(N\) = number of periods.

\subsubsection{SHADOW4 Undulator Models}

\textbf{A. "Undulator Gaussian" (Simplified)}
\begin{itemize}
    \item Assumes monochromatic source at/near resonance
    \item Rays sampled from Gaussian distributions
    \item Includes optional energy spread correction
    \item Useful for quick beamline prototyping
\end{itemize}

\textbf{B. "Undulator Light Source" (Full Model)}
\begin{enumerate}
    \item Compute far-field electric field stack (3D: energy, \(\theta\), \(\varphi\)) for filament beam
    \item Apply energy spread correction to angular array (if monochromatic)
    \item Use intensity stack as PDF to sample ray energies and directions
    \item Add electron emittance via Gaussian convolution
    \item Backpropagate far field to source plane for spatial distribution
    \item Sample ray positions from convolved distribution
    \item Assign polarization from interpolated far-field data
\end{enumerate}

\subsubsection{Important for 4th-Generation Sources}

\begin{itemize}
    \item \textbf{Diffraction-limited size} becomes significant with low electron emittance
    \item \textbf{Energy spread effects} are crucial for high harmonics
    \item \textbf{Non-Gaussian phase space} requires careful modeling beyond Gaussian approximations
    \item \textbf{Backpropagation methods} are essential for accurate source size calculation
\end{itemize}

% \subsubsection{Key Findings}

% \begin{itemize}
%     \item Ray tracing validates analytical models for size/divergence
%     \item Energy spread significantly broadens divergence at high harmonics
%     \item Off-resonance intensity profiles change shape (double-peak, etc.)
%     \item Polychromatic simulations match accumulated monochromatic results
% \end{itemize}

\subsubsection{Conclusion}

Undulator radiation is highly collimated and peaked at harmonic energies. Accurate ray-tracing requires proper treatment of electron beam parameters, non-Gaussian phase space modeling, backpropagation for source size calculation, and energy spread corrections for high harmonics. SHADOW4 implements both simplified Gaussian and full wave-optics-based models for modern synchrotron beamline design.

% \end{document}